% Post-CP18 source: energy application proved; known correspondence and its scope verified by four independent file audits.
\subsection{Energy-controlled reduction cycles and exact pointwise realization}

\begin{proposition}[Proved: reduction cycles with the established linear energy]
\label{D:energy-controlled-reduction-cycle}
Let $A$ be a complete excellent normal local surface $k$-algebra,
where $k$ is an algebraically closed coefficient and residue field,
and let $\pi:Y\to\operatorname{Spec}A$ be its minimal projective
regular resolution. Suppose its exceptional fibre is nonempty.
Write $\kappa_i=K_Y\cdot E_i\ge0$.
Choose an effective exceptional Cartier divisor $Z>0$ of full
support satisfying
\[
 -Z\cdot E_i\ge3\kappa_i\quad(\kappa_i>0),\qquad
 -Z\cdot E_i\ge0\quad(\kappa_i=0).
                                                        \tag{D.KR.1}
\]
Then $Z$, and every positive integral multiple of $Z$, is a
reduction cycle in the following precise sense:
\[
 \mathcal O_Z(-Z)\text{ is globally generated},\qquad
 H^1(Z,\mathcal O_Z(-Z))=0,\qquad
 \omega_Z^{-1}\text{ is globally generated}.
                                                        \tag{D.KR.2}
\]
In the cover family of Corollary~\ref{D:punctured-cover-linear-energy},
one may choose this $Z$ with the same bounds
\[
 -Z^2+K_Y\cdot Z+\sum_i(-Z\cdot E_i)=O(n),\qquad
 h^0(tZ,\mathcal O_{tZ})+h^1(tZ,\mathcal O_{tZ})=O_t(n)
                                                        \tag{D.KR.3}
\]
for each fixed $t\ge1$. Here $n$ is the cover degree and the
constants depend only on the fixed base and $t$. The assertion
does not enlarge that previously established cover family.
\end{proposition}

\begin{proof}
For a nonzero effective exceptional cycle $B=\sum_i m_iE_i$,
put $a=K_Y\cdot B=\sum_i\kappa_i m_i$. Adjunction gives
\[
                    2p_a(B)=2+a+B^2.                 \tag{D.KR.4}
\]
If $a\ge1$, negative definiteness implies
$2p_a(B)\le a+1\le2a$. If $a=0$, every component of its support
is a smooth rational $(-2)$-curve. The integral intersection
form on that support is even and negative definite, whence
$B^2\le-2$ and $p_a(B)\le0$.

For every subcycle $0<B\le tZ$, (D.KR.1) therefore gives
\[
 \begin{split}
 \deg_B\mathcal O_{tZ}(-tZ)&=-tZ\cdot B\ge2p_a(B),\\
 \deg_B\omega_{tZ}^{-1}&=-(K_Y+tZ)\cdot B
                              \ge2p_a(B).
 \end{split}                                           \tag{D.KR.5}
\]
In the second inequality, when $a>0$ the left side is at least
$(3t-1)a\ge2a$; when $a=0$ it is nonnegative.
The point-ideal automatic-adjunction argument in
Theorem~\ref{D:reflexive-pullback-uniform-generation}, applied
on the full possibly nonreduced curve $tZ$, now proves global
generation of both line bundles and $H^1$-vanishing for the
first. Equivalently, use the complete proof of
\citet[Lemma~2.4 and Theorem~1.1, pp.~192--193]{C_CFHR1999}.
This proves (D.KR.2), also for every $tZ$.

Choose the lattice parameter $c\ge5$ in
Lemma~\ref{D:integral-antinef-energy-cycle}; its lower bound is
stronger than (D.KR.1). Equations (D.KR.3) are then exactly the
bounds already proved in
Corollary~\ref{D:punctured-cover-antinef-cycle}. If a whole
connected exceptional graph has $\kappa=0$, use its rational
double point fundamental cycle, whose energy is two by
Proposition~\ref{D:RDP-full-sheaf-degree}. The same parity
argument proves (D.KR.5) there. Disjoint connected components
can be combined. Thus no extra denominator or enlargement of
the energy budget is needed to obtain a reduction cycle.
\end{proof}

\begin{theorem}[Known correspondence: full-module realization and determination on a reduction cycle]
\label{D:reduction-cycle-full-module-correspondence}
Let $A$ be a complete excellent normal local surface $k$-algebra
with algebraically closed coefficient and residue field $k$,
of arbitrary characteristic. Let $\pi:Y\to\operatorname{Spec}A$
be a projective regular resolution, an isomorphism away from
the closed point, and let $Z>0$ have full exceptional support.
Assume that $\mathcal O_Z(-Z)$ and $\omega_Z^{-1}$ are
generically generated by their global sections and that
$H^1(Z,\mathcal O_Z(-Z))=0$.

For a finite reflexive $A$-module $M$, put
\[
       \mathcal F_M=(\pi^*M)^{\vee\vee},\qquad
       R_Z(M)=\mathcal F_M|_Z.
\]
The functor $R_Z$ is injective on isomorphism classes. A vector
bundle $V$ on $Z$ is in its essential image if and only if:
\begin{enumerate}
\item $V$ is generically generated by global sections;
\item there is a vector bundle $V_2$ on $2Z$, with
$V_2|_Z\simeq V$, for which the boundary map
\[
 \delta_0:H^0(Z,V(Z))\longrightarrow H^1(Z,V)
                                                        \tag{D.KR.6}
\]
is injective. This is the boundary of
$0\to V\to V_2(Z)\to V(Z)\to0$.
\end{enumerate}
In addition, for a full bundle $\mathcal F=\mathcal F_M$
and $V=\mathcal F|_Z$, restriction induces
\[
 H^1(Y,\mathcal End\mathcal F)
       \simeq H^1(Z,\mathcal End V),
 \quad\text{equivalently}\quad
 \operatorname{Ext}^1_Y(\mathcal F,\mathcal F)
       \simeq\operatorname{Ext}^1_Z(V,V).
                                                        \tag{D.KR.7}
\]

The object correspondence is the known Kahn--Gavran theorem.
The proof below records its precise complete-ring scope and the
characteristic-independent marking and tangent-space verification;
no new classification theorem is claimed.
These are pointwise object and tangent-space statements.
They do not assert fullness on morphisms, lifting of a
prescribed isomorphism on $Z$, or an equivalence for arbitrary
Artinian parameter families.
\end{theorem}

\begin{proof}
We reconstruct the relevant argument in the complete algebraic
setting. The original theorem is
\citet[Theorem~1.4, Propositions~1.6 and~1.9, Claim on p.~148,
and Corollary~1.10, pp.~143--149]{C_Kahn1989}.
Its stated conventions include characteristic zero.
\citet[Section~2, Theorem~2 and Propositions~3--4,
pp.~69--71]{C_Gavran2012} extends the correspondence to complete
normal surface rings without that characteristic hypothesis,
but refers back to parts of Kahn's proof. The arguments below
include those parts and use no characteristic-zero vanishing.

\emph{Generic generation and the first cohomology.}
Suppose $\mathcal F$ is a vector bundle on $Y$ whose restriction
$V$ to $Z$ is generically generated. Products of generic
generators of $V$ and $\mathcal O_Z(-Z)$ give a map
\[
      \mathcal O_Z(-Z)^a\longrightarrow V(-jZ)
\]
with zero-dimensional cokernel, for every $j\ge1$.
Its kernel has no $H^2$, and its cokernel has no $H^1$.
Consequently $H^1(Z,V(-jZ))=0$. Filtering the restrictions of
$\mathcal F(-Z)$ to $nZ$ by these sheaves gives
$H^1(nZ,\mathcal F(-Z)|_{nZ})=0$ for every $n$.
Formal functions and completeness give
\[
 H^1(Y,\mathcal F(-Z))=0,
 \qquad H^1(Y,\mathcal F)\simeq H^1(Z,V).              \tag{D.KR.8}
\]
In particular all sections of $V$ lift to $Y$, so
$\mathcal F$ is generated at the exceptional generic points.
Away from the exceptional fibre the evaluation map from the
coherent module $\pi_*\mathcal F$ is onto because $\pi$ is an
isomorphism. This proves generic generation on $Y$.
The same reasoning, applied to $\mathcal O_Y(-jZ)$, supplies
global maps $\mathcal O_Y^a\to\mathcal O_Y(-jZ)$ that are
onto outside finitely many points.

\emph{A finite boundary detects fullness.}
For a vector bundle generically generated on $Y$, fullness is
equivalent to the surjectivity
\[
       H^0(Y,\mathcal F)\longrightarrow
       H^0(Y\setminus Z,\mathcal F).                  \tag{D.KR.9}
\]
Indeed the left module is finite over $A$. If (D.KR.9) holds,
Hartogs makes it reflexive; generic generation identifies its
reflexive pullback with $\mathcal F$ at all codimension-one
points and hence everywhere. Conversely a reflexive module
has exactly this Hartogs property.

Write $F_n=\mathcal F|_{nZ}$. Since $Z$ is Cartier and
$\mathcal F$ is locally free, the local-cohomology computation
for a single nonzerodivisor gives
\[
 H^1_Z(Y,\mathcal F)
        =\varinjlim_n H^0(nZ,F_n(nZ)).                 \tag{D.KR.10}
\]
This follows by taking sections of
$\mathcal F(*Z)/\mathcal F
 =\varinjlim_n F_n(nZ)$; the zeroth local-cohomology sheaf
vanishes. Thus (D.KR.9) is equivalent to injectivity of
$H^1_Z(Y,\mathcal F)\to H^1(Y,\mathcal F)$.
The $n=1$ map, followed by (D.KR.8), is exactly $\delta_0$.

If $\delta_0$ is injective, every map
\[
 \delta_j:H^0(Z,V((j+1)Z))\longrightarrow H^1(Z,V(jZ))
                                                        \tag{D.KR.11}
\]
obtained by twisting the $2Z$ sequence is injective as well.
To see this, dualize a generically surjective global map
$\mathcal O_Y^a\to\mathcal O_Y(-jZ)$. The resulting injection
$\mathcal O_Y(jZ)\to\mathcal O_Y^a$ gives a comparison of
(D.KR.11) with $a$ copies of (D.KR.6); the map on $H^0$
is injective. The comparison exists on $2Z$, so its boundary
maps commute. This proves the assertion.
The successive quotients in (D.KR.10) have boundary maps
whose compositions with restriction are (D.KR.11).
Their injectivity shows that the direct system of $H^0$ in
(D.KR.10) is already constant at $n=1$. We obtain the claimed
injection of local cohomology. The reverse implication follows
by restricting that injection to $n=1$.
Together with (D.KR.8), this proves the full-bundle criterion
using only $V$ and its actual $2Z$ extension.

\emph{Existence of a full extension.}
Start with $V_2$ satisfying the two stated conditions. At every
step, the obstruction to lifting a vector bundle on $nZ$ to
$(n+1)Z$ belongs to
$H^2(Z,\mathcal End V(-nZ))=0$. Choose successive lifts.
The resulting formal vector bundle algebraizes on $Y$ by
Grothendieck existence: the $Z$-adic and closed-fibre adic
topologies are cofinal, and the base is complete. The
algebraized coherent sheaf is locally free because its
nonfree locus is closed and proper over the local base and
misses the closed fibre. The preceding finite-boundary
criterion makes this bundle full. These formal-existence
inputs have the same precise proper complete-base scope as
\citep[Tags 0885, 08MF, 08MD and 02OC]{C_StacksBoundedReflexive}.

\emph{The boundary surjectivity needed for uniqueness.}
Fix a full extension $\mathcal F$ of $V$. For $j\ge1$, the
endomorphism sequence associated to the passage from $jZ$ to
$(j+1)Z$ has a boundary whose restriction is
\[
 \partial'_j:
 H^0(Z,\mathcal End V(-(j-1)Z))
       \longrightarrow H^1(Z,\mathcal End V(-jZ)).     \tag{D.KR.12}
\]
This map is surjective. Here are the details of the relevant
duality argument. The restrictions of
$\mathcal F$, $\mathcal O_Y(-(j-1)Z)$ and
$\omega_Y^{-1}(-Z)$ are generically generated on $Z$.
The last assertion is precisely the hypothesis on
$\omega_Z^{-1}$, by adjunction. The first paragraph lifts
their generic generators to $Y$. Their products therefore
give a generically surjective map
\[
 \mathcal O_Y^a\longrightarrow
                  \mathcal F(-jZ)\otimes\omega_Y^{-1}.
                                                        \tag{D.KR.13}
\]
Serre duality identifies the dual of the injective map
$\delta_0$ with the surjective boundary
\[
 H^0(Z,V^\vee\otimes\omega_Z)
      \longrightarrow H^1(Z,V^\vee(-Z)\otimes\omega_Z).
                                                        \tag{D.KR.14}
\]
Multiplication by (D.KR.13) compares $a$ copies of
(D.KR.14) with (D.KR.12). The map between the target $H^1$
groups is onto: its sheaf map has zero-dimensional cokernel,
and its kernel has no $H^2$. Both comparisons come from maps
on $2Z$, so the boundary square commutes. Surjectivity of
(D.KR.14) proves surjectivity of (D.KR.12). This is the
entire Claim on p.~148 of Kahn in the present algebraic scope.

\emph{Uniqueness of the full bundle.}
For a fixed bundle $F_j$ on $jZ$, isomorphism classes of its
locally free lifts to $(j+1)Z$, with a marking on $jZ$, form
a torsor under $H^1(Z,\mathcal End V(-jZ))$.
Changing the marking gives the actual $\operatorname{Aut}(F_j)$
action on this torsor.
For $j\ge2$, an element of
$H^0(Z,\mathcal End V(-(j-1)Z))$ defines an endomorphism
$h$ of $F_j$ with square zero, so $1+h$ is invertible.
It is the identity on $(j-1)Z$. In local transition matrices,
changing the marking by $1+h$ changes the lifting cocycle by
the boundary $\partial'_j(h)$. There are no higher terms,
because the relevant products of ideal powers vanish.
By (D.KR.12), every lifting class is reached.

At $j=1$, use the actual change-of-marking action, rather
than a pullback-only formula that would forget the induced
identification of the kernel. The group
$\operatorname{Aut}(V)$ is the group of units in the
finite-dimensional algebra $H^0(Z,\mathcal End V)$; it is
a smooth open subscheme of that vector space. The differential
of its orbit map at the identity is the boundary of the
endomorphism sequence, namely $\partial'_1$. This follows
by lifting $1+\epsilon h$ on an affine cover and computing
the discrepancy of the lifted transition matrices.
Thus the orbit map is smooth at the identity and its orbit
is open in the affine torsor of lifts. The full locus is
open there, because it is the injectivity locus of (D.KR.6)
between two fixed finite-dimensional vector spaces.
Every full lift has an open orbit by the same argument.
Two nonempty open subsets of this irreducible affine torsor
intersect, so there is exactly one full orbit.

The corrections for $j\ge2$ are the identity one level
earlier. Therefore the successive isomorphisms stabilize on
each fixed thickening and yield a compatible formal
isomorphism. Formal existence for coherent Hom algebraizes
it and its inverse. This proves uniqueness of the full
bundle, hence injectivity of $R_Z$ on module isomorphism
classes by $M=H^0(Y,\mathcal F_M)$.

\emph{The tangent-space comparison.}
The restriction sequence for $\mathcal End F_{j+1}$ has
kernel $\mathcal End V(-jZ)$. Its $H^0$ boundary contains
the surjective map (D.KR.12), while its kernel has no $H^2$.
It follows that
\[
 H^1((j+1)Z,\mathcal End F_{j+1})
             \simeq H^1(jZ,\mathcal End F_j)
 \quad(j\ge1).
\]
Formal functions gives (D.KR.7). All steps use curve duality,
generic generation, formal existence, or elementary
deformation theory and work in the stated characteristic.
\end{proof}

\begin{remark}[Exact scope of the applied finite-cycle information]
Proposition~\ref{D:energy-controlled-reduction-cycle} puts the
pointwise correspondence on the already constructed cycles
with degree-linear energy. In particular the conditions on a
bundle and its $2Z$ extension are exact pointwise realization
conditions, not merely necessary numerical conditions.
The fixed-determinant charts of
Proposition~\ref{D:hereditary-fixed-determinant-cycle-chart}
can be used on these cycles and on $2Z$ for the periodic
objects covered by that proposition.
Neither (D.KR.7) nor uniqueness on isomorphism classes says
that an arbitrarily prescribed map on $Z$ lifts: the exact
sequence still permits
\[
 H^1(Y,\mathcal End\mathcal F(-Z))
   \simeq\operatorname{coker}\bigl(
 H^0(Y,\mathcal End\mathcal F)\to
 H^0(Z,\mathcal End V)\bigr).
\]
Arbitrary-parameter realization, effective arrow equations,
and the desired degree-linear hereditary representation
count remain separate questions. No finite-generation or
nonabelian-crown conclusion is inferred here.
\end{remark}
